\subsection*{P2b: smaller integral seeds and the exact Weil lattice}
The polynomial projector is useful on arbitrary family members. At the specific elliptic product one can do substantially better by using explicit divisor classes. The following construction also fixes an integral normalization of the target plane.

\begin{proposition}[P2b: integral divisor-product seeds]\label{P2b}
For the balanced product $E^n\times E^n$, $n=2,4$, the integral Weil lattice has a basis $u_n,v_n$ with positive Gram matrix
\begin{equation}\label{eq:WeilGram}
 G_n=3^{n-1}\begin{pmatrix}2&1\\1&2\end{pmatrix}.
\end{equation}
Both basis vectors have integral signed cycle representatives. Their effective positive and negative degrees are bounded respectively by
\[
 (D_{u,2},D_{v,2})=(80,80),\qquad
 (D_{u,4},D_{v,4})=(147456,175104).
\]
Thus the raw Hecke transport of this integral basis has multiplier $m^{2n}$ and degrees at most $D_{u,n}m^{2n},D_{v,n}m^{2n}$. No projector denominator is needed for these particular seed cycles.
\end{proposition}
\begin{proof}
Write $e_j,f_j$ for the integral basis dual to the lattice $1,\omega$ on each elliptic factor, with $\int e_jf_j=1$. On a paired surface $E_x\times E_y$ set
\[
 (e_x+\omega f_x)(e_y+\omega^2 f_y)=U+\omega V.
\]
Exterior products are understood. Expansion gives
\[
 U=e_xe_y-e_xf_y+f_xf_y,\qquad V=-e_xf_y+f_xe_y.
\]
Let $F_x,F_y$ be the coordinate elliptic divisors and $\Gamma_1,\Gamma_\omega$ the two graphs. Their classes satisfy
\[
 V=[\Gamma_1]-[F_x]-[F_y],\qquad
 U=[\Gamma_\omega]+[\Gamma_1]-2[F_x]-2[F_y].
\]
For verification, the graph classes are respectively $\theta+V$ and $\theta+U-V$. Their pairings with the integral exterior basis agree with restriction to the maps $z\mapsto(z,z)$ and $z\mapsto(z,\omega z)$, which determines their Poincar\'e dual classes. Their self-intersections are zero, their mutual intersection is $\Nm(1-\omega)=3$, and their $\theta$-degrees are two. The coordinate divisors each have degree one. Accordingly $U$ has both sign degrees four and $V$ both sign degrees two.

On the $n$ paired surfaces define $u_n,v_n$ by
\[
 \zeta=\bigwedge_{i=1}^n(e_{x_i}+\omega f_{x_i})
        \wedge\bigwedge_{i=1}^n(e_{y_i}+\omega^2f_{y_i})
       =u_n+\omega v_n.
\]
The coefficients are integral. The all-$e$ coefficient is one; the coefficient obtained by replacing only $e_{x_1}$ by $f_{x_1}$ is $\omega$. Thus a rational combination of $u_n,v_n$ is integral only when both coefficients are integers. Since $\zeta,\bar\zeta$ span the two extreme $K$-eigenspaces, $u_n,v_n$ are exactly a basis of the integral Weil lattice.

One has $\zeta^2=\bar\zeta^2=0$ and $\int\zeta\bar\zeta=(-1)^n3^n$. Expanding these identities gives
$\int u_n^2=\int v_n^2=2(-1)^n3^{n-1}$ and $\int u_nv_n=(-1)^n3^{n-1}$, proving \eqref{eq:WeilGram}.

Up to the harmless overall sign $(-1)^{n(n-1)/2}$, $\zeta$ is the product of the $U_i+\omega V_i$. A term with $k$ factors $V$ and $n-k$ factors $U$ has each effective sign degree
\[
 n!\,2^{3n-k-1}.
\]
Indeed the total absolute degree on each paired factor is respectively eight or four, the positive and negative degrees are equal, and the product polarization contributes the multinomial factor $n!$. Write $\omega^k=a_k+b_k\omega$. The two bounds are therefore
\[
 D_{u,n}=n!\sum_{k=0}^n\binom nk|a_k|2^{3n-k-1},\qquad
 D_{v,n}=n!\sum_{k=0}^n\binom nk|b_k|2^{3n-k-1}.
\]
Using $\omega^0=1$, $\omega^1=\omega$, $\omega^2=-1-\omega$, $\omega^3=1$ and $\omega^4=\omega$ gives the four stated integers. Applying the same special-unitary isogeny calculation as in P2 proves the final assertion. The exterior algebra, lattice minors, Gram matrix, and degree sums are checked by \texttt{integral\_seeds.py}.
\end{proof}

For a fixed target $t=b_1u_n+b_2v_n$, choose $q_t>0$ with $q_tb_i\in\Z$. The corresponding explicit orbit formulas become
\[
 N_x=q_tm^{2n},\qquad
 D_x\le m^{2n}\bigl(|q_tb_1|D_{u,n}+|q_tb_2|D_{v,n}\bigr).
\]
They still describe chosen representatives, not minimal complexity. Dividing these cycles by $N_x$ gives rational cycles of bounded weighted degree on the whole orbit, while their integral clearing multipliers can grow. This explains why rational weighted degree cannot replace the integral-effective condition in R20.

For the principal datum $V=1$ and \eqref{eq:WeilGram} makes the denominator envelope in P3 especially concrete:
\[
 R=\max\left(1,\left\lfloor
 \frac{|v_a|\delta_n^2D^2}{2\,3^{n-1}(n!)^2}
 \right\rfloor\right),\qquad
 N=d_t\operatorname{lcm}(1,\ldots,R).
\]
Thus a numerical uniform seed degree would immediately produce a numerical denominator for any specified integral-basis normalization.
